\cotefile{156-7}
\scanpage{18}\mspage{17}

\diagram{\begin{tikzcd}[ampersand replacement=\&,row sep=large,column sep=small]
\operatorname{Omb} F \arrow[dr] \& \& \operatorname{Omb} \unsure{C|G} \arrow[dl] \\
\& \operatorname{Omb}(F)\cap\operatorname{Omb}(C) \arrow[d] \& \\
\& \operatorname{Omb}(L) \&
\end{tikzcd}}

\uline{At\unsure{spac|spec}} (Critère de \uline{compatibilité pour} ateliers \uline{dits spacieux}) Comme Atpol, mais en supprimant la condition de la fin ($L=\varnothing$ ou $L\in M$).

\uline{Proposition} \unsure{correspondante} pour ateliers spacieux : Comme \unsure{prop|prop.} ci-dessus, en supprimant \unsure{(vi)|(Vi)}.

% slanted
\margin{il suffit $=$ ici d'avoir \uline{\(\mathrm{At}\unsure{D'_1|D'}\)}, forme affaiblie de AtD, où on \unsure{affirme} \deleted{\unsure{supp.}} que $X'_L=\varnothing,Y'_L=\varnothing\Rightarrow X'\unsure{\mathrel{\Vert}|)(}Y'$ (NB $\mathrm{At}D'\unsure{\Rightarrow|\text{---}}\mathrm{At}D_0$)}
\diagram{\begin{tikzcd}[ampersand replacement=\&,row sep=0.15em,column sep=0.85em,execute at end picture={
\coordinate (junction) at ([xshift=1.4em,yshift=0.65em]nb.north east);
\draw (upper.west) .. controls +(left:0.6em) and +(right:1em) .. (junction);
\draw (lower.west) .. controls +(150:1.2em) and +(right:1em) .. (junction);
\draw[->] (junction) .. controls +(left:0.6em) and +(right:0.6em) .. ([xshift=-0.8em,yshift=0.35em]nb.north west);
}]
\& |[alias=upper]| \mathrm{At}\,D+\mathrm{At\,spac} \arrow[dr,Rightarrow] \& \& \& \\
|[alias=nb]| \text{\uline{NB}} \& \& \mathrm{At}\,C \arrow[r,Rightarrow] \& \mathrm{At}\,D \arrow[r,Rightarrow] \& \mathrm{At}\,D_0 \\
\& |[alias=lower]| \mathrm{At}\,D+\mathrm{At\,pol} \arrow[ur,Rightarrow] \& \& \&
\end{tikzcd}}

\uline{At fil} (Critère des Sup filtrants \inserted{loc. finis}) Soient $F\in\mathscr{F}$, $(F_i)_{i\in I}$ une famille \deleted{\unsure{filtrante}} croissante \inserted{(pour $\leq$)} de raffinements $F_i\ll F$, ``\unsure{localmt} finie \unsure{dans} $F$'' \unsure{i.e.} $\forall x\in\widetilde F$, la famille des $F_i\midX$ est stationnaire. Alors $\operatorname{Sup}F_i$ existe dans $\mathscr{F},\leq$.
% slanted
\margin{\unsure{Ou va-t-on} seulement supposer $i\leq j\Rightarrow F_i\ll\unsure{j|F_j}$ (ald $F_i\leq F_j$) ? (et prendre Sup pour $\ll$ ??)}

\uline{Corollaire 1} \deleted{\unsure{Rec}} (Moyennant At C) Recollement des raffinements, quand $F=\operatorname{Sup}F_i$ (deux versions, suivant qu'on recolle sur des sous-figures de $\underline F$, ou sur ses strates)
\scanpage{19}\mspage{18}

\uline{Corollaire 2} \unsure{Enoncé|Énoncé} analogue à Atfil, mais avec $(F_i)$ filtrants croissants pour $\ll$ ald $\leq$.

\uline{At supp} (Axiome \deleted{\unsure{du}}des supports d'un
\par\noindent\vrule\hspace{.5em}\begin{minipage}{.94\linewidth}
raffinement) Soi\deleted{\unsure{en}}\unsure{t|ent} $\deleted{F,}\deleted{\unsure{G|C}}X\in\deleted{\mathscr{F}}M$, \deleted{\unsure{avec}} $G\ll X$. Supposons $\exists Y\in M$ avec $Y\mathrel{\Vert}G$, $Y\mathrel{\overline{\Vert}}X$. Alors \deleted{on peut} il existe \unsure{une|un} telle $Y$ avec $Y\ll X$.
\end{minipage}

% slanted
\margin{Peut-être y-a-t-il lieu de renforcer : Pour $Y$ \unsure{donné}, $Y\mathrel{\Vert}G$, $Y\mathrel{\overline{\Vert}}X$, $\exists Y'\in M$, $Y'\ll F$, telle que $Y'\mathrel{\Vert}G$, $Y'\mathrel{\overline{\Vert}}Y$ \deleted{(au moins \unsure{supp.} $Y'$ \unsure{connexe|conn.} \unsure{?!|!})} Ceci est lié à la question de l'étude de l'atelier ``complémentaire'' de $G$\ldots\ On y reviendra\ldots}

\uline{Prop.} Idem que At supp, avec $X$ remplacé par figure $F$ quelconque ($G\ll F$)

Question : Quand $Y\mathrel{\Vert}G$, $Y\mathrel{\overline{\Vert}}F$, peut-on trouver $=$ $Y'\ll Y$ avec $Y'\ll F$ ? (On aura \unsure{automatiquement} $\deleted{\unsure{A}}Y'\mathrel{\Vert}G$, par At D.)

Ce \deleted{\unsure{est en|est}}\inserted{serait en} somme ça, le ``bon axiome''.

Mais (moyennant At L 1, \unsure{ancien|voire} At 8) on est ramené au cas où $Y$ est un lieu, \unsure{soit} $x$, et : se pose la question
\[
 x\mathrel{\overline{\Vert}}Y \Longleftrightarrow
 \underset{\inserted{\text{i.e. }x\in\operatorname{omb}(Y)}}{x\ll Y}
 \qquad (\Longleftarrow\ \text{est évident}
\]
i.e.
\[
 x\mathrel{\Vert}Y \Longleftrightarrow
 \deleted{x\mathrel{\Vert}\unsure{\operatorname{om}|\operatorname{omb}}}\;x\notin\operatorname{omb}(Y)
\]

Mais ceci n'est guère vérifié que pour les ateliers \unsure{ensemblistes}. Notons pourtant

\uline{At ens}\ \uline{D\deleted{\unsure{ens}}}. \inserted{(\uline{Critère} de \uline{disjonction ensembliste})} \deleted{Pour que} Soient $X,Y\in M$. Pour que $X\mathrel{\Vert}Y$, il \deleted{suff} (f. et) suffit que $\deleted{\unsure{Z}}\operatorname{Omb}(\unsure{x|X})\cap\operatorname{Omb}(Y)=\varnothing$ i.e. qu'il \uline{n'}existe pas de $Z\in M$ avec $Z\ll X,Z\ll Y$
\scanpage{20}\mspage{19}

\uline{Cor 1} Idem pour figures $F,\unsure{C|G}$ :
\[
 F\mathrel{\Vert}C\Longleftrightarrow \operatorname{Omb}F\cap\operatorname{Omb}C=\varnothing
\]

\uline{Cor 2} \deleted{Soit}\unsure{Soient} $x\in\mathcal{L}$, $F\in\mathscr{F}$. Alors
\[
 x\mathrel{\Vert}F\Longleftrightarrow x\not\ll F\quad\text{i.e.}\quad x\notin\operatorname{omb}(F)
\]

\uline{Cor. 3} Soient $x,y\in\mathcal{L}$, alors
\[
 x\mathrel{\Vert}y\Longleftrightarrow x\ne y\quad .
\]

\uline{NB} On a trivialement
\[
 \begin{aligned}
 &\deleted{\mathrm{At\,supp}}\\
 &\unsure{\mathrm{Ad}|\mathrm{At}}\,\mathrm{ens}\,D\quad\Longrightarrow\quad\mathrm{At\,supp}.
 \end{aligned}
\]
On notera que Atsupp est vérifié dans tous les ateliers auxquels j'ai songé aux fins pratiques, mais non At ens D.

\begin{center}\rule{.24\linewidth}{.4pt}\end{center}

J'en viens aux axiomes impliquant les lieux (At L 1, 2, 3)

\noindent\vrule\hspace{.5em}\begin{minipage}{.94\linewidth}
\uline{At L 1} \inserted{(\uline{Axiome de divisibilité})} Soit $X\in M$, alors $\operatorname{omb}(X)^{\unsure{\circ|o}}\ne\varnothing$ i.e. $\exists x\in\mathcal{L}$ avec $x\mathrel{\overset{\circ}{\ll}}X$.
\end{minipage}

Cet axiome d'anodine apparence implique \unsure{pourtant} (et signifie, \unsure{essentiellt|essentiellement}) que les ``espaces'' ou ``\unsure{endroits|antres}'' où vivent nos ateliers, sont ``infiniment divisibles'', et par suite (\deleted{\unsure{ils}}\unsure{quand} ils ne sont discrets) sont infinis : c'est l'introduction de ``l'infini actuel'' \deleted{\unsure{tacit|tout}} !
\scanpage{21}\mspage{20}

\uline{At L 2} (Critère de \uline{disjonction par les lieux}) Soient $X,Y\in M$. Alors
\[
 \begin{aligned}
 X\mathrel{\Vert}Y\quad\Longleftrightarrow\quad&
 \deleted{\unsure{\exists|\forall}}\ \forall x\in\operatorname{omb}(\unsure{Y|X}),\
 y\in\operatorname{omb}Y,\ \text{on a}\\
 &\hspace{4em} x\mathrel{\Vert}y
 \end{aligned}
\]
\margin{(À ne pas confondre avec $\operatorname{omb}(X)\cap\operatorname{omb}(Y)=\varnothing$ !)}

\margin{\uline{Cor} Soient $F,\unsure{C|G}\in\mathscr{F}$. Alors \(F\mathrel{\Vert}C\Longleftrightarrow\ \forall\substack{x\in\operatorname{omb}(\unsure{F|X})\\y\in\operatorname{omb}(C)}\ \text{on a }x\mathrel{\Vert}y\)}

\uline{At L 3} (\deleted{\unsure{Crit}}Axiome de disjonction pour les lieux) Soient $X,Y\in M$, $\deleted{\unsure{X\ne Y}}\,X\mathrel{>\mkern-10mu<}Y$, $X\ne Y$. Alors
\[
 x\in\operatorname{omb}(X)^{\unsure{\circ|o}},\quad
 y\in\operatorname{omb}(Y)^\circ\quad\Longrightarrow\quad x\mathrel{\Vert}y
\]
% slanted
\margin{équivaut, moyennant L2,}

\unsure{Equivaut|Équivaut} à

\uline{Cor} Soit $F\in\mathscr{F}$. Alors deux lieux appartenant à des strates ouvertes différentes de $F$ sont disjoints.

\uline{NB} Quand on admet \uline{At} D, alors il suffit de poser At L 3 dans le cas $Y\leq X$.

Si on admet At L 2, alors At C(\deleted{\unsure{pol}}spac) prend la forme

\noindent\vrule\hspace{.5em}\begin{minipage}{.94\linewidth}
\uline{At CL} (\deleted{pol}spac) (Critère de compatibilité par les lieux, pour ateliers \deleted{\unsure{polyédraux|polyédriques}}\inserted{spacieux}) Soient $X,Y\in M$, $L=X\cap Y=\underset{Z\in\widetilde X\cap\widetilde Y}{\operatorname{Sup}}Z$. Alors pour que $X$ et $Y$ soient \unsure{comp.|compa.}, \deleted{\unsure{il f.|il faut}}\uline{il suffit}\textsuperscript{*} $\forall x\in\operatorname{omb}(X)\setminus\operatorname{omb}(L)$, $y\in\operatorname{omb}(Y)-\operatorname{omb}(L)$, on a $x\mathrel{\Vert}y$.
\end{minipage}
% slanted
\margin{* Ou vaut-il mieux poser : il faut et \unsure{(?)|?} suffit ?}

En effet, \deleted{la \unsure{condit}} dans l'énoncé de \uline{At spac} avec $X,Y,L$ : $X'\ll X$, $Y'\ll Y$, la condition $X'_L=\varnothing,Y'_L=\varnothing\Rightarrow X'\mathrel{\Vert}Y'$, cette dernière condition s'explicite, par L2, par
\scanpage{22}\mspage{21}
\begingroup\csname tagsleft@true\endcsname

\unsure{termes de|lemme de} $x\in\operatorname{omb}(X')\subset\operatorname{omb}(X\unsure{.})$, $y\in\operatorname{om}(Y')\subset\operatorname{omb}(Y)$ par $x\mathrel{\Vert}y$. Or la condition $X'_L=\varnothing$ implique $x_L=\varnothing$ car $x_L\ll X'_L$, et \deleted{\unsure{ce|elle}}\unsure{cela} signifie \deleted{que} ici que $\unsure{x|x'}\notin\operatorname{omb}(L)$ (cf lemme) et idem pour $y$.

Donc $x\in\operatorname{omb}(X)\setminus\operatorname{omb}(L)$, $y\in\operatorname{omb}(Y)\setminus\operatorname{omb}(L)$ et si on a At CL(spac), on conclut bien $x\mathrel{\Vert}y$, \deleted{\unsure{d'où}} Donc la condition
\[
 \forall x\in\operatorname{omb}(X)\setminus\operatorname{omb}(\deleted{\unsure{X|Y}}L),\
 y\in\operatorname{omb}(Y)\setminus\operatorname{omb}(L),\
 \text{\deleted{on} a}\quad x\mathrel{\Vert}y
 \tag*{(AL)}
\]
implique
\[
 \forall \unsure{X|X'}\in\operatorname{Omb}(\unsure{x|X|x'}),\
 Y'\in\operatorname{Omb}(Y),\ X'_L=Y'_L=\varnothing,\ \text{on a}\quad X\mathrel{\Vert}Y
 \tag*{(A)}
\]
et l'inverse est trivial, comme on voit en prenant $X,Y\in\mathcal{L}$. On a utilisé ici le

\uline{Lemme 1} \unsure{Soit|Soient} $\deleted{\unsure{L|E}\leq F}\ P\leq Q$ (en l'occurrence $\deleted{\unsure{X}}L\leq X$ ou $L\leq Y$) $x\in\unsure{\operatorname{omb}|\operatorname{Omb}}(Q)$, alors $\deleted{x\mathrel{\Vert}Q\unsure{\Rightarrow|-}}$
\[
 x_P=\varnothing\Longleftrightarrow x\notin\operatorname{omb}(\deleted{\unsure{Q}}P)\quad
 \text{i.e}\quad x\not\ll\deleted{\unsure{Q}}P\quad
 (\deleted{\unsure{\Longleftrightarrow}x\mathrel{\Vert}\unsure{P|Q}})
\]

En effet, $\deleted{x\ll P\Longleftrightarrow x_P\unsure{\Rightarrow|=x}}$ pour la première équivalence, passant aux négations on \deleted{\unsure{trouve}}\inserted{doit} prouver $x\ll P\Longleftrightarrow x_P\ne\varnothing_{\unsure{\mathscr{F}|F|7}}$, et comme $x_P\ll x$ et $x$ minimal dans $M\setminus\{\varnothing_{\mathscr{F}}\}$ pour $\ll$, $x_P\ne\varnothing_{\mathscr{F}}$ signifie aussi $x_P=x$ i.e. $x\ll P$ OK.

\deleted{D'autre part, $x\mathrel{\Vert}P\Rightarrow x\not\ll P$ est trivial, \unsure{vu} la définition} Je profite de l'occasion pour expliciter

\noindent\vrule\hspace{.5em}\begin{minipage}{.94\linewidth}
\uline{Lemme 2} \inserted{(moyennant At D)} Soient $P\leq Q$, et $F\ll Q$. Conditions équivalentes
\[
 F_P=\varnothing\Longleftrightarrow F\mathrel{\Vert}P
\]
\end{minipage}

On applique At D (ou plutôt, la prop. qui suit, \unsure{haut|haut de} page \unsure{15|16}) à $P,Q$ $P\ll P$, $F\ll Q$, pour avoir $|\Rightarrow$ ; l'autre est vrai avec les seuls axiomes At 1 -- At 7
\endgroup
\scanpage{23}\mspage{22}

On introduit de $=$

\uline{At CL} (p\uline{ol)} (Critère de compatibilité par les lieux, pour ateliers polyédraux)

Notons

\usetikzlibrary{decorations.pathreplacing}
\diagram{\resizebox{\linewidth}{!}{\begin{tikzcd}[ampersand replacement=\&,row sep=0.9em,column sep=0.8em,
execute at end picture={
\draw[decorate,decoration={brace,mirror,amplitude=3pt}] ([xshift=4pt]rowa.north east) -- ([xshift=4pt]rowe.south east);
\draw (crossleft.west) -- (crossright.east) -- (oldzero.east);
}]
\& |[alias=crossleft]| \deleted{\underline{\mathrm{At}\,D_0}+\underline{\mathrm{At}\,L1}+\underline{\mathrm{At}\,CL}(\mathrm{spac})} \arrow[rr,Leftrightarrow]
\& \& \deleted{\underline{\mathrm{At}}\,L2} \arrow[r,Rightarrow]
\& |[alias=crossright]| \deleted{\mathrm{At}\,D} \arrow[d,Rightarrow] \\
|[alias=rowa]| \text{a)} \&
\deleted{\mathrm{At}\,\unsure{L|L1}}\ \mathrm{At}\,L2+\mathrm{At}\,CL
\underset{\substack{\deleted{\unsure{\text{ou}}}\\(\text{ou pol})}}{(\mathrm{spac})}
\arrow[rr,Leftrightarrow] \& \&
\mathrm{At}\,L2+\mathrm{At}\,C\underset{(\text{ou pol})}{(\mathrm{spac})}.
\& |[alias=oldzero]| \deleted{\mathrm{At}\unsure{C}D_0} \\
\text{b)} \&
\underset{\deleted{\unsure{\mathrm{At}\,D_0}+}}{\inserted{\mathrm{At}\,D_0+}}\,
\underline{\mathrm{At}\,L}1+\underline{\mathrm{At}\,CL}\underset{(\text{ou pol})}{(\mathrm{spac})}
\arrow[rr,Rightarrow] \& \& \underline{\mathrm{At}\,L2}\arrow[d,Rightarrow] \& \deleted{\unsure{\Rightarrow}\mathrm{At}\,D} \\
\text{c)} \& \underline{\mathrm{At}\,L}2+\underline{\mathrm{At}\,L3} \arrow[r,Rightarrow] \&
\underline{\mathrm{At}\,D}\arrow[r,Rightarrow] \& \underline{\mathrm{At}}\,D_0 \& \\
\text{d)} \& \underline{\mathrm{At}\,L}3+\text{\uuline{At CL}}(\mathrm{spac}) \arrow[rr,Rightarrow] \& \& \underline{\mathrm{At}\,C} \arrow[ul,Rightarrow]
\& \deleted{\unsure{(\Rightarrow\mathrm{At}\,D\Rightarrow\mathrm{At}\,D_0)}}\\
|[alias=rowe]| \text{e)} \& \underline{\mathrm{At}\,D}\arrow[r,Rightarrow] \&
\deleted{\unsure{\mathrm{At}}}L3 \& \& \\
\& \deleted{\unsure{\mathrm{At}\,L1\ \text{ou }A}}\&\&\&
\end{tikzcd}}}
% slanted
\margin{Voir \unsure{diagr\textsuperscript{me}|diagramme} refait p. 45}

Pratiquement, je retiens ceci, concernant l'ensemble des axiomes de compatibilité-disjonction [et de divisibilité]
\diagram{\begin{tikzcd}[ampersand replacement=\&,row sep=0.4em,column sep=0.65em]
C\arrow[r,Rightarrow]\&D\arrow[r,Rightarrow]\arrow[d,Rightarrow]\&D_0,\&
\mathrm{At}\,\deleted{\unsure{C}}C(\mathrm{spac})\&\mathrm{At}\,CL(\mathrm{spac})\\
{[L1]}\ L2,\&L3\&\&\mathrm{At}\,C(\mathrm{pol})\&\mathrm{At}\,CL(\mathrm{pol})
\end{tikzcd}}

1) Du point \deleted{\unsure{des}} vue des axiomes ``non conditionnels'' (sans \inserted{mention} spac ou pol \ldots), les deux axiomes L2, L3 de disjonction impliquant les lieux, impliquent D donc $D_0$, mais non C (\deleted{\unsure{axiome|axiomes} \unsure{ne|n'} \unsure{pas}} de compatibilité, ne se \unsure{ramenant|réduisant} pas à une disjonction) Pour déduire C d'axiomes pertinents sur les lieux, il nous faut des At CL spac ou At CL (pol). À ce sujet, on note

$\underline{\mathrm{At}\,L}2+\underline{\mathrm{At}\,L}3+\mathrm{At}\,CL(\mathrm{spac})$ implique tous les axiomes de compatibilité-disjonction (\unsure{de|c'est} style \unsure{spac.})) 
\diagram{\(\begin{array}{ll}
\mathrm{At}\,C(\mathrm{spac}) & \text{par (a)}\\
\mathrm{At}\,D & \text{par (c) d'où }D_0\\
\mathrm{At}\,C & \text{par NB p. 17 ou par d}
\end{array}\)}

On peut \unsure{remplacer} \deleted{\unsure{donc que|dire que} At L} ``pol'' ald ``spac''
\scanpage{24}\mspage{23}

L'axiome de divisibilité At L1 est à part, il n'est pas inclus dans ce paquet, mais on a

\diagram{\resizebox{\linewidth}{!}{\begin{tikzcd}[ampersand replacement=\&,row sep=0.55em,column sep=1.1em,execute at end picture={
\draw[->] ([xshift=-1.6em,yshift=-2pt]shorta.west) -- (shorta.west);
\draw[->] ([xshift=-1.6em,yshift=2pt]shortb.west) -- (shortb.west);
}]
\mathrm{At}\,L1+\mathrm{At}\,L\deleted{\unsure{1}}2\underset{\deleted{+\mathrm{At}\,D_0}}{+}\mathrm{At}\,L3+\mathrm{At}\,CL(\mathrm{spac})
\arrow[d,Leftrightarrow]\&
\text{\parbox{13em}{impliquent tous les axiomes de compatibilité ou de support, \unsure{i.e.|int.} style ``spac'' en plus bien sûr At L1 lui-$=$.}}\\
\mathrm{At}\,D_0+\mathrm{At}\,L1+\mathrm{At}\,L3+\mathrm{At}\,CL(\mathrm{spac})\arrow[d,Rightarrow]\&\\
|[alias=shorta]| \mathrm{At}\,D+\mathrm{At}\,L\unsure{1|3}+\mathrm{At}\,CL(\mathrm{spac})\arrow[d,Rightarrow]\&\\
|[alias=shortb]| \mathrm{At}\,D+\mathrm{At}\,L1+\mathrm{At}\,L2+\mathrm{At}\,C(\mathrm{spac})\&
\end{tikzcd}}}
% slanted
\margin{\deleted{\unsure{ici, il y a}} deux axiomes ``absolus'' At D, At C(sp) et deux axiomes relatifs à L}
et on \unsure{peut remplacer|peut écrire} ``pol'' ald ``spac''.

Récapitulation des axiomes (à l'exclusion des \deleted{\unsure{d|de}} axiomes d'exactitude d'une subdivision, version générale, ou ``spac'', ou ``pol''\ldots)

% margin note written sideways
\margin{axiomes ``sine qua non''}
% slanted
\margin{comparaison de $\leq$ et $\ll$}
\usetikzlibrary{decorations.pathreplacing}
\diagram{\resizebox{\linewidth}{!}{\begin{tikzcd}[ampersand replacement=\&,row sep=0.32em,column sep=1em,cells={nodes={anchor=west}},
execute at end picture={
\draw[decorate,decoration={brace,mirror,amplitude=4pt}] ([xshift=-3pt]aone.north west) -- ([xshift=-3pt]athree.south west) node[midway,left=5pt] {$\leq$};
\draw[decorate,decoration={brace,mirror,amplitude=4pt}] ([xshift=-3pt]afour.north west) -- ([xshift=-3pt]aseven.south west);
\draw[decorate,decoration={brace,mirror,amplitude=5pt}] ([xshift=-3em]aone.north west) -- ([xshift=-3em]aseven.south west);
}]
|[alias=aone]| \mathrm{At}\,1\&\text{Sup des familles filtrantes majorées}\\
\mathrm{At}\,2\&\text{Critère de }\leq\text{ par strates}\\
|[alias=athree]| \mathrm{At}\,3\&\text{Critère de \deleted{\unsure{disjon|disj}}\unsure{compatibilité|discompatibilité} par strates}\\
|[alias=afour]| \mathrm{At}\,4\&\text{Critère de }\ll\text{ par strates}\\
\mathrm{At}\,5\&\text{Critère du raffinement entre sous-figures }\ (\text{Si }F,G\leq H,\ F\leq G\Leftrightarrow F\ll G)\\
\mathrm{At}\,6\&\text{Axiome de la plus petite strate, ou des \unsure{raffint|raffimt} \unsure{internes|intérieurs} de strates }\ X\overset{\unsure{\circ|o}}{\ll}Y\\
|[alias=aseven]| \mathrm{At}\,7\&\text{Transitivité du raffinement interne}
\end{tikzcd}}}

% slanted
\margin{axiomes de compatibilité-disjonction}
\diagram{\resizebox{\linewidth}{!}{\begin{tikzcd}[ampersand replacement=\&,row sep=0.24em,column sep=1em,cells={nodes={anchor=west}},
execute at end picture={
\draw[decorate,decoration={brace,mirror,amplitude=4pt}] ([xshift=-4pt]comp.north west) -- ([xshift=-4pt]poly.south west);
\draw[rounded corners=4pt] ([xshift=-4pt,yshift=3pt]fil.north west) rectangle ([xshift=4pt,yshift=-3pt]fend.east |- sup.south);
\draw[->] ([xshift=4pt]fend.east) to[out=-10,in=44] (spend.north east);
\draw[->] ([xshift=-3em]polend.south east) -- ++(-0.3em,-1.1em);
\draw[decorate,decoration={brace,mirror,amplitude=5pt}] ([xshift=-3pt]lone.north west) -- ([xshift=-3pt]lpol.south west);
}]
|[alias=comp]| \mathrm{At}\,C\arrow[d,Rightarrow,start anchor=south west,end anchor=north west]\&
\text{Critère de compatibilité de raffinement de sous-figures.}\\
\mathrm{At}\,D\arrow[d,Rightarrow,start anchor=south west,end anchor=north west]\&
\text{Critère de disjonction de}\qquad\text{- - - -}\\
\mathrm{At}\,D_0\&\text{\deleted{disjonction} stabilité de }\Vert\text{ par }\ll\\
\mathrm{At}\,C(\mathrm{spac})\&\text{Critère de compatibilité dans les ateliers spacieux}\\
|[alias=poly]| \mathrm{At}\,C(\mathrm{pol})\&\text{- - - -}\qquad\qquad\text{- - -}\quad\text{polyédraux}\\
|[alias=fil]| \mathrm{At}\,\text{Fil.}\&|[alias=fend]| \text{Existence de Sup filtrants de}\inserted{\text{une famille loc. finie de}}\text{ raffinements}\\
|[alias=sup]| \mathrm{At}\,\text{supp.}\&\ldots\\
|[alias=lone]| \mathrm{At}\,L1\&\text{Axiome de divisibilité}\\
\mathrm{At}\,L2\&\text{Critère de disjonction par les lieux}\\
\deleted{\unsure{\ll}}\,\mathrm{At}\,L3\&\text{Axiome de disjonction pour les lieux}\\
\mathrm{At}\,CL(\mathrm{spac})\&|[alias=spend]| \text{Critère de compatibilité par les lieux (at. spacieux)}\\
|[alias=lpol]| \mathrm{At}\,CL(\mathrm{pol})\&|[alias=polend]| \rule{14em}{0.32pt}\ \text{(-- polyédraux)}
\end{tikzcd}}}
% slanted
\margin{axiomes concernant les lieux}
\figure{\begin{tikzpicture}
\path[use as bounding box] (0,0) rectangle (.96\linewidth,.14);
\fill (0.35,.01) -- (.40,.12) -- (.45,.13) -- (.46,.10) -- (.42,.09) -- (.39,0) -- cycle;
\end{tikzpicture}}
\scanpage{25}\mspage{24}
\begingroup\csname tagsleft@true\endcsname

Je reviens sur la description \unsure{d'un|d'une} atelier spacieux ou polyédral, satisfaisant L1, L2, L3 plus \deleted{At \unsure{L(spac)|L1(spac)} ou} At CL(spac) (ou, au choix At CL(pol)), en termes de $(M,\leq,\ll,\mathrel{>\mkern-10mu<})^{*}$.
% slanted
\margin{* ce qui détermine \unsure{tout} au moins les \deleted{\unsure{familles}} figures de \unsure{réf.} \deleted{\unsure{ou|et}} \unsure{et} les sous-figures de telles\ldots}

Le point ici, c'est que le critère At CL(spac) réduit, pour connaître \({\mathrel{>\mkern-10mu<}}_M\), à la connaissance de $\deleted{\mathrel{>\mkern-10mu<}}{\Vert}_{\mathcal L}$, i.e. à celle d'une structure de graphe sur $\mathcal L$. (NB ds le cas d'un atelier ensembliste, je rappelle que $x\mathrel{\Vert}y\Longleftrightarrow x\neq y$)

Donc, la question est, partant de \deleted{$(M,\leq,\ll)$}, \deleted{\unsure{d'où|de}, \unsure{et} $\mathcal L$.}
\[
\unsure{,}\quad(M,\leq,\ll)\tag*{(1)}
\]
d'où
\[
\mathcal L=\{x\in M\mid x\text{ minimal pour }\ll\},\tag*{(2)}
\]
et d'une relation symétrique antiréflexive
\[
x\mathrel{\Vert}y\qquad(x,y\in\mathcal L)
\]
(\unsure{impliquant} $x\neq y$) sur $\mathcal L$, quelles conditions imposer pour que ça corresponde bien à un atelier\textsuperscript{\#} spacieux, satisfaisant \deleted{\unsure{At}} (en plus At \deleted{\unsure{L}}CL(spac)) les conditions L1, L2, L3.
% slanted
\margin{\# Il est entendu qu'on prendra pour $\unsure{F|\mathscr F}$, soit l'ens. de \uline{toutes} les \deleted{familles}, parties fermées de $M_{\leq}$ formées d'\unsure{él\textsuperscript{ts}|él\textsuperscript{t}} deux à deux compatibles, soit \unsure{qu'on précise celles|qu'on précise les} qui sont des s.f. (en supposant que tt partie fermée d'un $\widetilde X$ en est \unsure{une s.f.|de s.f.}).}

Il est entendu qu'on va \uline{définir} \({\mathrel{>\mkern-10mu<}}_M\) par
\[
\begin{aligned}
X\unsure{,|{}_{1}}\mathrel{>\mkern-10mu<}Y\quad&\Longleftrightarrow\quad
\forall\,\begin{aligned}[t]&x\in\operatorname{omb}(X)\setminus\operatorname{omb}(L)\\
&y\in\operatorname{omb}(Y)-\operatorname{omb}(L)\end{aligned},\\
&\hspace{5em}\text{on a}\quad x\mathrel{\Vert}y
\end{aligned}\tag*{(3)}
\]
\margin{\diagram{\resizebox{\linewidth}{!}{\begin{tikzcd}[ampersand replacement=\&,row sep=0.25em,cells={nodes={anchor=west}},
execute at end picture={
\draw ([xshift=-3pt]where.south west) .. controls +(right:.4em) and +(up:.8em) .. ([xshift=-1pt]bound.west) .. controls +(down:.8em) and +(left:.6em) .. (bound.south east);
}]
|[alias=where]| \text{où}\quad L=X\cap Y\\
\phantom{\text{où}\quad L}=\underset{Z\in\widetilde X\cap\widetilde Y}{\operatorname{Sup}}Z\\
\unsure{\text{donc}|\text{d'où}|\text{dans ce}}\\
|[alias=bound]| \operatorname{omb}(L)\overset{\text{déf}}{=}\underset{Z\in\widetilde X\cap\widetilde Y}{\bigcup}\operatorname{omb}(Z\unsure{)}
\end{tikzcd}}}}
Pour commencer, il faut \deleted{\unsure{vérifier}} \inserted{écrire} la condition préliminaire
\[
X\mathrel{>\mkern-10mu<}Y,\ X'\leq X,\ Y'\leq Y\ \Longrightarrow\ X'\mathrel{>\mkern-10mu<}Y'\tag*{(4)}
\]
\deleted{et on a besoin ce faisant \unsure{de|en} sous \unsure{formes|forme} plus}
\endgroup
